% ============================================================
% SECTION 4.7: MECHANISM. Written 2026-07-26 (late) after the audit.
%
% SCOPE, established by the audit and easy to overstate: this explains the
% DEGREE-BASED channel, not the boundary. The boundary's mechanism is the
% denoising evidence in 4.6. PAPER_STRUCTURE.md called the old section
% "Mechanism: why the boundary is there", which was wrong.
%
% Numbers, all recomputed this session:
%   exp_B/results.csv        17 networks; dQ_fixed positive 17/17 real and
%                            17/17 null, null larger in 15; delta larger on the
%                            null in 16/17, median ratio 1.599
%   exp_M/probe_results*.csv decomposition on the NULL rows (the terms are
%                            defined relative to the real arm, so real rows are
%                            zero by construction): sorting positive 16/16,
%                            median +0.5743, range +0.08 to +3.42; identity
%                            term_sum = log_ratio to 1.7e-14; facebook-combined
%                            excluded, its real delta is -0.0014
%   exp_Q/verdicts.csv       pfix arm only: delta moves 37 to 1322 neutral sd,
%                            sign change 5/5; re-detection AMI 0.824-0.944 on the
%                            UP ends and 0.410-0.540 on the DOWN ends; the
%                            literal arm is excluded because it destroys the
%                            partition (AMI <= 0.038, p_intra -> 0)
% STYLE: no em-dashes, no AI-characteristic phrasing (see STYLE.md).
% ============================================================

\subsection{Mechanism of the degree-based channel}\label{sec:experiments:mechanism}

A degree-based sparsifier keeps an edge with probability proportional to a score
computed from its endpoints' degrees. Whether that helps or harms a partition
depends on one quantity: the difference between the mean score on edges inside
communities and the mean score on edges between them. Write it $\delta$. When
$\delta > 0$ the sampler removes between-community edges preferentially, which
raises the intra-community fraction of any fixed partition and therefore its
modularity. This subsection establishes what controls $\delta$, and it concerns
the degree-based channel only. The mechanism behind the boundary of
Section~\ref{sec:experiments:boundary} is the denoising evidence given there.

\paragraph{$\delta$ is not evidence about communities}
Two quantities in this literature are read as showing that a sparsifier clarifies
community structure: $\delta$ itself, and the modularity change of a fixed
partition when the graph is sparsified. Both were measured on seventeen real
networks ranging from $10^3$ to $3.8 \times 10^6$ vertices, and on
degree-preserving rewirings of each, built by double edge swaps, which hold every
vertex degree exactly and destroy the communities.

Neither survives. The fixed-partition modularity change is positive on all
seventeen real graphs; on their rewirings it is positive on all seventeen as
well, and larger than the real value on fifteen. The separation $\delta$ is
larger on the rewiring than on the real graph in sixteen of seventeen cases, with
a median ratio of $1.60$. Real community structure suppresses the separation
rather than producing it. Neither quantity is used as evidence anywhere else in
this paper, and the fixed-partition measure in particular is the same
self-scoring comparison that Section~\ref{sec:protocol:transfer} rejects.

\paragraph{Where the null's advantage comes from}
The ratio between the null's separation and the real graph's decomposes exactly.
Writing $\delta = r_{pb}\, \sigma_s / \sqrt{p(1-p)}$, with $r_{pb}$ the
point-biserial correlation between an edge's score and whether it lies inside a
community, $\sigma_s$ the spread of the scores and $p$ the intra-community edge
fraction, gives
\begin{equation}\label{eq:delta-decomp}
\log \frac{\delta_{\mathrm{null}}}{\delta_{\mathrm{real}}}
= \underbrace{\log \frac{r_{pb}^{\mathrm{null}}}{r_{pb}^{\mathrm{real}}}}_{\text{sorting}}
+ \underbrace{\log \frac{\sigma_s^{\mathrm{null}}}{\sigma_s^{\mathrm{real}}}}_{\text{scale}}
+ \underbrace{\tfrac{1}{2}\log \frac{p^{\mathrm{real}}(1-p^{\mathrm{real}})}
{p^{\mathrm{null}}(1-p^{\mathrm{null}})}}_{\text{granularity}} ,
\end{equation}
an identity that holds in our measurements to $1.7 \times 10^{-14}$. The sorting
term accounts for the difference. It is positive on all sixteen networks where
the ratio is defined, with a median of $+0.57$ and a range from $+0.08$ to
$+3.42$, while the scale and granularity terms are small and mostly negative. The
seventeenth network is excluded because its real separation is $-0.0014$, which
makes the ratio undefined. What the rewiring changes is not the spread of the
scores or the number of communities but how well the score sorts within-community
edges from between-community ones, and it sorts them better once the communities
are gone.

\paragraph{What controls $\delta$}
The correlational evidence points at where high-degree edges sit relative to
community boundaries, but a graph statistic that co-varies with $\delta$ across
networks does not establish that moving it moves $\delta$. We therefore steer it
directly. Starting from a real graph and its detected partition, edges are
rewired to increase or decrease the concentration of high-degree-product edges on
community boundaries, while the degree sequence, the partition, the
intra-community edge fraction and the modularity are all held exactly fixed. Only
the placement of hub edges relative to the boundaries changes.

$\delta$ follows. Across five networks it moves between 37 and 1322 standard
deviations of the unsteered chains, in the direction it is steered, and it
changes sign in all five when steered downward. The largest excursion takes
ca-GrQc from $+0.087$ to $-0.875$ with its degree sequence, partition,
intra-edge fraction and modularity unchanged.

Two limits on that result are worth stating. The steered graphs remain
recognizable in the upward direction, where re-detecting communities from scratch
recovers the original partition at AMI $0.82$ to $0.94$, and less so downward,
where re-detection agrees at $0.41$ to $0.54$. And a looser variant of the same
procedure, which does not hold the intra-community fraction fixed, drives that
fraction to zero and destroys the partition altogether, with re-detection AMI
below $0.04$; nothing about the mechanism can be concluded from that variant and
we exclude it.

\paragraph{An explanation that fails}
Triangle structure is the natural alternative: edges inside communities close more
triangles, so a score correlated with triangle participation would produce
$\delta > 0$ without any reference to degree placement. The steering experiment
rules it out. Intra-community triangle density falls in both directions of the
steering, from $11.7$ to $9.0$ upward and to $11.4$ downward on ca-GrQc and
similarly elsewhere, while $\delta$ moves in opposite directions in the two arms.
A quantity that decreases while the target rises in one arm and falls in the
other does not control the target.

\paragraph{What this explains}
The degree-based channel is mechanical. It is governed by where high-degree-product
edges sit relative to community boundaries, a property that a degree-preserving
rewiring can reproduce and often exaggerate, and it is therefore not a route by
which a degree-based sparsifier clarifies community structure. That is consistent
with the negative territory of Section~\ref{sec:experiments:negative}, where the
degree-based arm behaves like the others, and with the finding in
Section~\ref{sec:experiments:boundary} that above the transition the gain does not
depend on which local scoring rule is used.
